\section{Multi-reviewer adversarial critique}
\label{sec:reviewers}

The earlier self-critique (Section~\ref{sec:critique}) is one voice
arguing with the report. This section is five \emph{distinct}
reviewer voices, each grounded in a different professional bias, each
allowed to attack the work from their own angle. Each reviewer
delivers (i) the harshest \emph{honest} read of the result, (ii) a
specific evidence request, (iii) the response we can ship today from
the artefacts on disk, and (iv) what we cannot yet defend. The list
of reviewers is deliberately uncomfortable: a federated-systems
purist, a clinical biostatistician, an optimisation theorist, an
MLOps engineer, and an ICU intensivist. The point of the exercise is
not consensus; the point is to show the reader where each
constituency's complaint lands and how far the artefacts go to meet
it.

\subsection{Reviewer 1: Federated-systems purist}
\textbf{Likely complaint.} ``This is not a federated paper; it is an
ablation study using federated language. There is no
heterogeneity-aware aggregator beyond a discretised \cvar{}, no
secure aggregation, no client dropout, no asynchronous training, no
Byzantine analysis. The privacy-plus-bytes trade-off that motivates
the field is not measured.''

\textbf{Specific evidence request.} ``Show client-dropout robustness;
show Byzantine robustness with at least $1$ malicious client; show
that the bytes you report are the bytes a real deployment would
incur, not the upper bound at $1000$ rounds.''

\textbf{Response we can ship.} (i) The early-stopping ledger
(\texttt{stopped\_round} column in
\texttt{proposal\_method\_seed\_results.csv}) shows that no method
ran to the round budget; the median \fedprox{}-$0.1$ run terminated
at round $\sim$$25$, paying $\sim$$28\%$ of the round-budgeted bytes.
So the bytes we report are \emph{actual} early-stop bytes, not
worst-case. (ii) Per-seed early-stop variance is itself a partial
robustness result: the same operating point reaches the same
operating loss in $80\%$ of seeds within $\pm 5$ rounds of the
median. (iii) We do \emph{not} measure Byzantine robustness; that is
the next study.

\textbf{What we cannot defend yet.} Secure aggregation, client
dropout under realistic ICU outage profiles, asynchronous training,
gradient inversion attacks. These are squarely in the future-work
section.

\subsection{Reviewer 2: Clinical biostatistician}
\textbf{Likely complaint.} ``The operating-point recommendation is
made on average \auprc{}, but the deployment audience is a clinician
at one ICU. Per-client confidence intervals, calibration plots, and
decision-curve analysis are absent. The seed pool is $10$, which is
not enough to anchor a clinical trial.''

\textbf{Specific evidence request.} ``Per-client BCa-bootstrap
confidence intervals on \auprc{}, sensitivity at fixed specificity,
calibration slope and intercept per client, and decision-curve net
benefit at the deployment threshold.''

\textbf{Response we can ship.} (i) BCa CIs are computed and stored in
\texttt{outputs/full\_mimic\_iv\_proposal\_alignment/stats/proposal\_confidence\_intervals.csv}
($1{,}000$ resamples per CI) -- \fedprox{}-$0.1$ \auprc{} CI is
$[0.660, 0.670]$, non-overlapping with \fedavg{}'s $[0.624, 0.640]$.
(ii) Calibration is partly addressed by the ECE/MCE columns in the
clinical metrics CSV; a full reliability diagram is in
Figure~\ref{fig:per-client-clinical}. (iii) The seed pool is $10$;
that is small for biostatistics but it is large for FL ablation.

\textbf{What we cannot defend yet.} Decision-curve analysis at the
clinically meaningful threshold (e.g.\ $0.10$). Calibration
\emph{intercept} plots per client. Subgroup analysis by age, sex,
ethnicity. These are listed in Section~\ref{sec:limits}.

\subsection{Reviewer 3: Optimisation theorist}
\textbf{Likely complaint.} ``The convergence claims are empirical;
there is no rate analysis. The proximal term's effect on the
condition number of the local subproblem is not bounded. The LP
duality is interpreted as a shadow price, but the LP is over a
\emph{discrete} candidate set, so the dual is only valid on the
relaxation. The claim that DE finds an off-grid optimum needs a
local-curvature confirmation.''

\textbf{Specific evidence request.} ``Compute the local Hessian
spectrum at the best checkpoint; report the local condition number
of $F_k$ on each client; bound the gap between the discrete LP
optimum and its continuous relaxation; report DE's variance over
restarts.''

\textbf{Response we can ship.} (i) The 1D and 2D loss-landscape data
(\texttt{loss\_landscape/}) are the closest empirical analogue to
local curvature; the bowl widths are roughly the inverse of the
local Hessian's smallest eigenvalue. (ii) The KKT diagnostics
(\texttt{kkt\_diagnostics.csv}) confirm the LP is solved cleanly;
the discrete-vs-continuous gap is bounded by the spacing between
adjacent Pareto points in $\mathcal{X}$, which is reported in
\texttt{dense\_vs\_sparse\_shadow\_price.csv}. (iii) DE was not
restarted; the convergence trace shows plateau by iteration $9$, so
restarts are unlikely to change the operating point.

\textbf{What we cannot defend yet.} A formal convergence rate for
\fedprox{} on \mimic{}; a tight bound on the LP's
discreteness-induced gap; a Hessian-spectrum analysis. These are
appendix-quality follow-ups.

\subsection{Reviewer 4: MLOps engineer}
\textbf{Likely complaint.} ``The deployment story is incomplete:
where is the inference latency, the model size on disk, the cold-
start cost on a fresh ICU, the canary deployment plan, the rollback
criterion? The watchdog is reported as never triggering, which is
the only thing the deployment cares about, but the rest of the
operational picture is absent.''

\textbf{Specific evidence request.} ``Wall-clock per-round time on
the headline operating point; model size; expected per-inference
latency on a CPU; cold-start time when adding a $10$th client;
rollback criterion; SLO definition.''

\textbf{Response we can ship.} (i) Resource ledger is in
\texttt{outputs/full\_mimic\_iv\_proposal\_alignment/monitoring/resource\_summary.csv}:
peak memory $9.28$ GB, min free $41.75$ GB, no watchdog trigger.
(ii) Round time on M4 Max is $\sim$$5$ s for the MLP and
$\sim$$1$ s for the logistic; on production GPUs we expect both
to be sub-second. (iii) Model size: MLP $\sim$$26$ KB, logistic
$\sim$$248$ B (literally $30$ floats and a bias). (iv) The
cold-start story is partial: a new client can be added by
broadcasting the global model and running one local epoch
($\sim$$2$ s); we did not test the full handshake.

\textbf{What we cannot defend yet.} A canary deployment plan, an
explicit rollback criterion, and an SLO are squarely in the future-
work pipeline.

\subsection{Reviewer 5: ICU intensivist (the actual end-user)}
\textbf{Likely complaint.} ``Which patient does this model help, and
when does it kill someone? Show me the false-negatives at the small
units; show me the operating threshold; show me what happens to a
patient at admission, $24$ hours, $48$ hours; show me the model's
decision in the cases where the resident disagreed.''

\textbf{Specific evidence request.} ``A confusion matrix per ICU at a
clinically meaningful threshold; a per-patient timeline (the model is
single-shot but a deployment is repeated); a discussion of how the
model interacts with existing scoring (SOFA, APACHE).''

\textbf{Response we can ship.} (i) Per-client clinical metrics
(\texttt{per\_client\_clinical\_metrics.csv}) include sensitivity,
specificity, PPV, NPV at the $0.5$ threshold. (ii) The Neuro
Intermediate failure case is documented: $10$ deaths, $4$ caught,
$6$ missed -- those $6$ are the operationally important false
negatives. (iii) We do \emph{not} have a per-patient timeline; the
single-shot prediction is the proposal's stated scope.

\textbf{What we cannot defend yet.} A clinical-decision-support
integration plan, an interaction analysis with SOFA/APACHE, a
per-patient longitudinal model. These are explicitly future work.

\subsection{Synthesis: where the report is strong, where it is open}
\paragraph{Strong.} (i) The mechanism for every claim is grounded in
on-disk evidence; (ii) the LP duality is solved cleanly with KKT
diagnostics; (iii) the per-seed reliability is reported and is the
basis of the deployment recommendation; (iv) the watchdog never
triggered, so the resource budget is honest; (v) the dossiers are
auditable to the line of code.

\paragraph{Open.} (i) Byzantine robustness; (ii) decision-curve
analysis at the clinical threshold; (iii) local Hessian spectrum;
(iv) canary deployment plan; (v) per-patient longitudinal modelling.

The report does \emph{not} claim to be a finished paper for any of
the five reviewers' constituencies. It claims to be a faithful
mechanistic study of the federated-vs-centralized-vs-local
trade-off on \mimic{}, with the operational consequences explicitly
labelled.
